Mit $1\,\mathrm{PS} = \wPSO$ (FoTaBe S.~186) ist die mechanische Leistung $P = \PbzO = \Pbzw$. In der Zeit $t = \tbzO$ muss das Benzin wegen des Wirkungsgrads die Wärmeenergie
\begin{align}
 E &= \EbzF = \frac{\Pbzw\cdot\tbz}{\wgrO} = \Ebz
\end{align}
liefern. Mit dem Heizwert $H = \HbzO$ ist dazu die Benzinmasse
\begin{align}
 m &= \frac{E}{H} = \mbzF \\
   &= \frac{\Pbzw\cdot\tbz}{\wgrO\cdot\Hbz} \\
   &= \mbz \approx \result{\mbzP{0}{2}}
\end{align}
nötig. Mit der Dichte $\rho = \rhoBO$ (FoTaBe S.~187) entspricht das dem Volumen
\begin{align}
 V &= \VbzF = \frac{\mbz}{\rhoB} = \Vbz \approx \result{\VbzlP{0}{2}}
\end{align}
